\begin{abstract}
Este trabajo analiza la dinámica estocástica de la prima de riesgo griega respecto al bono alemán, utilizando datos diarios que abarcan el periodo crítico de la crisis de deuda soberana. La justificación de este estudio radica en la necesidad de contrastar metodologías dispares para modelar riesgos extremos: el enfoque físico de difusión anómala y el enfoque econométrico estándar. 

Nuestro objetivo es determinar si el mercado exhibe memoria a largo plazo o si se ajusta a los modelos de volatilidad condicional tradicionales. Para ello, utilizamos datos del diferencial de bonos a 10 años y aplicamos análisis de Desplazamiento Cuadrático Medio (MSPD) e inferencia Bayesiana mediante cadenas de Markov (MCMC). 

Los hallazgos empíricos revelan un comportamiento super-difusivo ($\alpha \approx 0.94$), evidenciando una fuerte persistencia estructural, mientras que el modelo GARCH demuestra ser indispensable para capturar clústeres de volatilidad a corto plazo, a pesar de su complejidad computacional.
\end{abstract}